% GENERATED FILE. Edit canonical lessons, then run npm run generate:books.
% canonical-source-hash: fnv1a64:b6b296534c8527af
% canonical-lessons: HI-C87-par, HI-C87-thaa, HI-C87-thii, HI-C87-the, HI-C87-nahin-thaa, HI-C87-bolta-thaa, HI-C87-repaso-atit, HI-C87-sintesis-atit

\chapter{The Past, And Where You Were}
\label{ch:hi-atit}
\hlchaptermodality{95}

\begin{chapteropening}
\textbf{By the end of this chapter:} \emph{I can say where I was and what I used to do, put the sentence in the negative, and use the postposition that pins a person to a place.}

Ask where somebody was and answer for yourself, using the exact sentence the first practice paper prints as reading item 14 and moving it one step into the past.
\end{chapteropening}

\section[par]{\hi{पर} — the word that follows the place instead of leading it}

\label{lesson:HI-C87-par}



\begin{quote}
You can name \textbf{\hi{घर}} and \textbf{\hi{स्टेशन}}. Try to say \emph{at home} — and notice that you have the place and no way to stand in it.
\end{quote}

\subsection*{You'll want to know: \hi{पर}}
\textbf{\hi{पर}} (\emph{par}) — \textquotedblleft{}at,\textquotedblright{} pinning you to one point.

\textbf{\hi{घर} \hi{पर}} — at home. \textbf{\hi{स्टेशन} \hi{पर}} — at the station.

\begin{cousinweb}
English sets \emph{at} in front of the place. Hindi sets \hi{पर} behind it:

\noindent
\begin{tabularx}{\linewidth}{@{}>{\raggedright\arraybackslash}X>{\raggedright\arraybackslash}X@{}}
\toprule
\textbf{English} & \textbf{Hindi} \\
\midrule
at home & \hi{घर} \hi{पर} \\
at the station & \hi{स्टेशन} \hi{पर} \\
\bottomrule
\end{tabularx}

That ordering is why these words are called \textbf{post}positions. Every one of them trails its noun, so once you have the noun you know exactly where the small word goes.

What \hi{पर} wants is a \textbf{point} rather than an interior — a house you are at, a platform you are standing on. A place you are \emph{inside} reaches for a different small word, which arrives later.
\end{cousinweb}

\begin{grammarlens}[title={Grammar Lens: where the pieces land}]
\noindent
\begin{tabularx}{\linewidth}{@{}>{\raggedright\arraybackslash}X>{\raggedright\arraybackslash}X@{}}
\toprule
\textbf{sentence} & \textbf{meaning} \\
\midrule
\hi{मैं} \hi{घर} \hi{पर} \hi{हूँ।} & I am at home. \\
\hi{वह} \hi{स्टेशन} \hi{पर} \hi{है।} & He is at the station. \\
\bottomrule
\end{tabularx}

Person, then place, then the small word, then the verb. The verb goes last and stays last, which is the shape every Hindi sentence you have met so far has had.
\end{grammarlens}

\subsection*{Your turn}
Say these aloud:

\begin{itemize}
\raggedright
  \item \textquotedblleft{}\hi{घर} \hi{पर}\textquotedblright{} — at home
  \item \textquotedblleft{}\hi{मैं} \hi{घर} \hi{पर} \hi{हूँ}\textquotedblright{}
  \item \textquotedblleft{}\hi{वह} \hi{स्टेशन} \hi{पर} \hi{है}\textquotedblright{}
  \item whether \hi{पर} goes before or after the place it marks
\end{itemize}

\begin{tcolorbox}[breakable,colback=teal!4,colframe=teal!35!black,title={Before you move on}]
Where does \hi{पर} sit, before the place or after it? (After.) Say \emph{at the station}. (\textbf{\hi{स्टेशन} \hi{पर}}.)
\end{tcolorbox}
\section[thā]{\hi{था} — one word swapped, and the sentence is behind you}

\label{lesson:HI-C87-thaa}



\begin{quote}
Say \emph{I am at home} in Hindi. Now try to say where you were an hour ago, and watch the sentence stop at its last word.
\end{quote}

\subsection*{You'll want to know: \hi{था}}
\textbf{\hi{था}} (\emph{thā}) — \textquotedblleft{}was,\textquotedblright{} spoken by or about a male subject.

Every other word in the sentence stays where it is. Take out the present-tense word, put \hi{था} in its place, and you have moved:

\noindent
\begin{tabularx}{\linewidth}{@{}>{\raggedright\arraybackslash}X>{\raggedright\arraybackslash}X@{}}
\toprule
\textbf{now} & \textbf{earlier} \\
\midrule
\hi{मैं} \hi{घर} \hi{पर} \hi{हूँ।} & \hi{मैं} \hi{घर} \hi{पर} \hi{था।} \\
\hi{वह} \hi{स्टेशन} \hi{पर} \hi{है।} & \hi{वह} \hi{स्टेशन} \hi{पर} \hi{था।} \\
\bottomrule
\end{tabularx}

\begin{cousinweb}
\hi{था} is worn down from Sanskrit \textbf{\hi{स्थित}} (\emph{sthita}), \textquotedblleft{}standing, situated.\textquotedblright{} To have been somewhere was, once, to have been standing there.

English arrives at the same picture along a different road. \emph{Station} comes from Latin \emph{stāre}, to stand — and Latin \emph{stāre} and Sanskrit \emph{sthā-} are the same inherited word, separated long before either language was written down. The Hindi for \emph{was} and the English for \emph{station} are distant relatives.
\end{cousinweb}

\begin{grammarlens}[title={Grammar Lens: two present words, one past word}]
\noindent
\begin{tabularx}{\linewidth}{@{}>{\raggedright\arraybackslash}X>{\raggedright\arraybackslash}X>{\raggedright\arraybackslash}X@{}}
\toprule
\textbf{subject} & \textbf{now} & \textbf{earlier} \\
\midrule
\hi{मैं} & \hi{हूँ} & \hi{था} \\
\hi{वह} & \hi{है} & \hi{था} \\
\bottomrule
\end{tabularx}

In the present, the verb changes to match \textbf{who} you are talking about. In the past it stops caring about that. One word covers both rows — which is a small mercy, and a short-lived one, because the next lesson shows what the past listens to instead.
\end{grammarlens}

\subsection*{Your turn}
\begin{itemize}
\raggedright
  \item \emph{Look:} at \textbf{\hi{था}} and name the character carrying its first sound (\textbf{\hi{थ}}), the one you met in \textbf{\hi{हाथ}}
  \item \emph{Say it:} \textquotedblleft{}\hi{था}\textquotedblright{} — \emph{thā}, was
  \item \emph{Say it:} \textquotedblleft{}\hi{मैं} \hi{घर} \hi{पर} \hi{था}\textquotedblright{}
  \item \emph{Say it:} \textquotedblleft{}\hi{वह} \hi{स्टेशन} \hi{पर} \hi{था}\textquotedblright{}
  \item \emph{Say it:} which word in \emph{\hi{मैं} \hi{घर} \hi{पर} \hi{हूँ}} you replace, and which words you leave alone
\end{itemize}

\begin{tcolorbox}[breakable,colback=teal!4,colframe=teal!35!black,title={Before you move on}]
How many words change when you move \emph{\hi{मैं} \hi{घर} \hi{पर} \hi{हूँ}} into the past? (One.) What did the Sanskrit ancestor of \hi{था} mean? (Standing, situated.)
\end{tcolorbox}
\section[thī]{\hi{थी} — the ending you have bent a dozen times, arriving somewhere new}

\label{lesson:HI-C87-thii}



\begin{quote}
\hi{था} ends in \textbf{-\hi{आ}}. Before reading on: what has every other word ending in \textbf{-\hi{आ}} done when a woman was the subject?
\end{quote}

\subsection*{You'll want to know: \hi{थी}}
\textbf{\hi{थी}} (\emph{thī}) — \textquotedblleft{}was,\textquotedblright{} spoken by or about a female subject.

\noindent
\begin{tabularx}{\linewidth}{@{}>{\raggedright\arraybackslash}X>{\raggedright\arraybackslash}X@{}}
\toprule
\textbf{speaker} & \textbf{earlier} \\
\midrule
a man & \hi{मैं} \hi{घर} \hi{पर} \hi{था।} \\
a woman & \hi{मैं} \hi{घर} \hi{पर} \hi{थी।} \\
\bottomrule
\end{tabularx}

\begin{cousinweb}
This is not a new pattern. It is the oldest one you have:

\noindent
\begin{tabularx}{\linewidth}{@{}>{\raggedright\arraybackslash}X>{\raggedright\arraybackslash}X@{}}
\toprule
\textbf{masculine} & \textbf{feminine} \\
\midrule
\hi{काला} & \hi{काली} \\
\hi{बोलता} & \hi{बोलती} \\
\hi{था} & \hi{थी} \\
\bottomrule
\end{tabularx}

A word ending in \textbf{-\hi{आ}} swaps that ending for \textbf{-\hi{ई}}. You learned it on adjectives, met it again on the habitual verb, and here it is a third time. The past copula behaves like an adjective because, historically, it once was one — a participle meaning \emph{situated}, agreeing with whoever was situated.
\end{cousinweb}

\begin{grammarlens}[title={Grammar Lens: what each tense listens to}]
\noindent
\begin{tabularx}{\linewidth}{@{}>{\raggedright\arraybackslash}X>{\raggedright\arraybackslash}X>{\raggedright\arraybackslash}X@{}}
\toprule
\textbf{tense} & \textbf{listens to the person} & \textbf{listens to gender} \\
\midrule
present & yes (\hi{हूँ}, \hi{है}, \hi{हैं}) & no \\
past & no & yes (\hi{था}, \hi{थी}) \\
\bottomrule
\end{tabularx}

The two tenses divide the work differently, and neither one asks you to track both at once. A woman says \hi{थी} whether she is talking about herself or about another woman.
\end{grammarlens}

\subsection*{Your turn}
Say these aloud:

\begin{itemize}
\raggedright
  \item \textquotedblleft{}\hi{थी}\textquotedblright{} — \emph{thī}, was
  \item \textquotedblleft{}\hi{मैं} \hi{घर} \hi{पर} \hi{थी}\textquotedblright{}
  \item \textquotedblleft{}\hi{वह} \hi{स्टेशन} \hi{पर} \hi{थी}\textquotedblright{}
  \item the two endings that mark the pair, and one adjective that shows them
\end{itemize}

\begin{tcolorbox}[breakable,colback=teal!4,colframe=teal!35!black,title={Before you move on}]
Which ending replaces \textbf{-\hi{आ}} for a female subject? (\textbf{-\hi{ई}}.) Does the present-tense verb bend for gender the way \hi{था} does? (No — that is what is new here.)
\end{tcolorbox}
\section[the]{\hi{थे} — the third ending, and the politeness you have already paid for}

\label{lesson:HI-C87-the}



\begin{quote}
You have \hi{था} and \hi{थी} for one person. Name the pronoun for \emph{we}, and guess which ending a group will want.
\end{quote}

\subsection*{You'll want to know: \hi{थे}}
\textbf{\hi{थे}} (\emph{the}) — \textquotedblleft{}were.\textquotedblright{}

\noindent
\begin{tabularx}{\linewidth}{@{}>{\raggedright\arraybackslash}X>{\raggedright\arraybackslash}X@{}}
\toprule
\textbf{subject} & \textbf{earlier} \\
\midrule
\hi{हम} & \hi{हम} \hi{घर} \hi{पर} \hi{थे।} \\
\hi{वे} & \hi{वे} \hi{स्टेशन} \hi{पर} \hi{थे।} \\
\hi{आप} & \hi{आप} \hi{कहाँ} \hi{थे}? \\
\bottomrule
\end{tabularx}

\begin{cousinweb}
\textbf{-\hi{ए}} is the plural ending on the same family of words as \textbf{-\hi{आ}} and \textbf{-\hi{ई}}, so \hi{थे} sits beside \hi{था} and \hi{थी} without being a new thing to learn.

The row worth pausing on is the last one. \textbf{\hi{आप}} takes \hi{थे} even when you are speaking to exactly one person, for the same reason it takes \hi{हैं} for one person: the polite level in Hindi is built out of plural forms. Asking a single guest \emph{\hi{आप} \hi{कहाँ} \hi{थे}?} is grammatically a question to a crowd, and that is the courtesy.

For a group of women, \hi{थे} adds a nasal and becomes \textbf{\hi{थीं}} — worth recognising on a page now, and not worth drilling yet.
\end{cousinweb}

\begin{grammarlens}[title={Grammar Lens: the level decision, once more}]
\noindent
\begin{tabularx}{\linewidth}{@{}>{\raggedright\arraybackslash}X>{\raggedright\arraybackslash}X>{\raggedright\arraybackslash}X@{}}
\toprule
\textbf{system} & \textbf{at the \hi{तुम} level} & \textbf{at the \hi{आप} level} \\
\midrule
copula now & \hi{हो} & \hi{हैं} \\
command & \hi{बोलो} & \hi{बोलिए} \\
copula earlier & \hi{थे} & \hi{थे} \\
\bottomrule
\end{tabularx}

The past does not split the two levels at all: \hi{तुम} and \hi{आप} both take \hi{थे}. One fewer decision, in the one place you might have expected another.
\end{grammarlens}

\subsection*{Your turn}
Say these aloud:

\begin{itemize}
\raggedright
  \item \textquotedblleft{}\hi{थे}\textquotedblright{} — \emph{the}, were
  \item \textquotedblleft{}\hi{हम} \hi{घर} \hi{पर} \hi{थे}\textquotedblright{}
  \item \textquotedblleft{}\hi{वे} \hi{स्टेशन} \hi{पर} \hi{थे}\textquotedblright{}
  \item the question \emph{where were you?} to one person you address as \hi{आप}
\end{itemize}

\begin{tcolorbox}[breakable,colback=teal!4,colframe=teal!35!black,title={Before you move on}]
Which ending does \hi{आप} take in the past, for one person? (\textbf{\hi{थे}}.) Why does a single person get a plural word? (The polite level is built out of plurals.)
\end{tcolorbox}
\section[nahīṃ thā]{\hi{नहीं} \hi{था} — saying where you were not}

\label{lesson:HI-C87-nahin-thaa}



\begin{quote}
Where does \textbf{\hi{नहीं}} stand in a present-tense sentence? Put it there in a past-tense one and check yourself against what follows.
\end{quote}

\subsection*{You'll want to know: \hi{नहीं} \hi{था}}
\textbf{\hi{नहीं} \hi{था}} (\emph{nahīṃ thā}) — \textquotedblleft{}was not.\textquotedblright{}

\hi{नहीं} keeps the seat it has always had: directly in front of the verb, at the end of the sentence.

\noindent
\begin{tabularx}{\linewidth}{@{}>{\raggedright\arraybackslash}X>{\raggedright\arraybackslash}X@{}}
\toprule
\textbf{present} & \textbf{past} \\
\midrule
\hi{मैं} \hi{घर} \hi{पर} \hi{नहीं} \hi{हूँ।} & \hi{मैं} \hi{घर} \hi{पर} \hi{नहीं} \hi{था।} \\
\hi{वह} \hi{स्टेशन} \hi{पर} \hi{नहीं} \hi{है।} & \hi{वह} \hi{स्टेशन} \hi{पर} \hi{नहीं} \hi{थी।} \\
\bottomrule
\end{tabularx}

\begin{cousinweb}
English rebuilds the sentence to negate it: \emph{I am at home} becomes \emph{I am \textbf{not} at home}, and \emph{I was} becomes \emph{I was not} — but \emph{I go} has to become \emph{I do not go}, with a helper verb appearing out of nowhere.

Hindi asks for none of that. One word, one seat, every tense. Learning the past costs you nothing extra here, because there was nothing extra to learn.
\end{cousinweb}

\begin{grammarlens}[title={Grammar Lens: the seat in front of the verb}]
\noindent
\begin{tabularx}{\linewidth}{@{}>{\raggedright\arraybackslash}X>{\raggedright\arraybackslash}X@{}}
\toprule
\textbf{sentence} & \textbf{who is speaking} \\
\midrule
\hi{मैं} \hi{घर} \hi{पर} \hi{नहीं} \hi{था।} & a man \\
\hi{मैं} \hi{घर} \hi{पर} \hi{नहीं} \hi{थी।} & a woman \\
\hi{हम} \hi{घर} \hi{पर} \hi{नहीं} \hi{थे।} & a group \\
\bottomrule
\end{tabularx}

\hi{नहीं} never changes shape. Everything that bends in these three lines bends after it, not because of it.
\end{grammarlens}

\subsection*{Your turn}
Say these aloud:

\begin{itemize}
\raggedright
  \item \textquotedblleft{}\hi{मैं} \hi{घर} \hi{पर} \hi{नहीं} \hi{था}\textquotedblright{}
  \item the same sentence as a woman would say it
  \item \textquotedblleft{}\hi{हम} \hi{स्टेशन} \hi{पर} \hi{नहीं} \hi{थे}\textquotedblright{}
  \item which word in the sentence \hi{नहीं} stands in front of
\end{itemize}

\begin{tcolorbox}[breakable,colback=teal!4,colframe=teal!35!black,title={Before you move on}]
Does \hi{नहीं} take a different shape in the past? (No.) Where does it stand? (Immediately before the verb, which is at the end.)
\end{tcolorbox}
\section[boltā thā]{\hi{बोलता} \hi{था} — what you used to do}

\label{lesson:HI-C87-bolta-thaa}



\begin{quote}
Say \emph{I speak Hindi}. Then look at the last word of it and ask what you would put there instead.
\end{quote}

\subsection*{You'll want to know: \hi{बोलता} \hi{था}}
\textbf{\hi{बोलता} \hi{था}} (\emph{boltā thā}) — \textquotedblleft{}used to speak.\textquotedblright{}

\noindent
\begin{tabularx}{\linewidth}{@{}>{\raggedright\arraybackslash}X>{\raggedright\arraybackslash}X@{}}
\toprule
\textbf{now} & \textbf{earlier} \\
\midrule
\hi{मैं} \hi{हिंदी} \hi{बोलता} \hi{हूँ।} & \hi{मैं} \hi{हिंदी} \hi{बोलता} \hi{था।} \\
\bottomrule
\end{tabularx}

One word replaced, and the habit is one you no longer have.

\begin{cousinweb}
Two pieces, two jobs:

\noindent
\begin{tabularx}{\linewidth}{@{}>{\raggedright\arraybackslash}X>{\raggedright\arraybackslash}X@{}}
\toprule
\textbf{piece} & \textbf{its job} \\
\midrule
\hi{बोलता} & this happens repeatedly \\
\hi{हूँ} / \hi{था} & and it happens \textbf{now} / it happened \textbf{then} \\
\bottomrule
\end{tabularx}

\textbf{-\hi{ता}} carries the habit. The word after it carries the time. They were never welded together, which is why moving the sentence into the past costs a single swap rather than a new conjugation.

English needs a whole phrase for this — \emph{used to} — and Hindi needs one letter's worth of change.
\end{cousinweb}

\begin{grammarlens}[title={Grammar Lens: both halves bend}]
\noindent
\begin{tabularx}{\linewidth}{@{}>{\raggedright\arraybackslash}X>{\raggedright\arraybackslash}X@{}}
\toprule
\textbf{speaker} & \textbf{earlier} \\
\midrule
a man & \hi{मैं} \hi{हिंदी} \hi{बोलता} \hi{था।} \\
a woman & \hi{मैं} \hi{हिंदी} \hi{बोलती} \hi{थी।} \\
\bottomrule
\end{tabularx}

Both words end in \textbf{-\hi{आ}} for a man and both swap to \textbf{-\hi{ई}} for a woman. They move together, so you are still tracking one decision rather than two.
\end{grammarlens}

\subsection*{Your turn}
Say these aloud:

\begin{itemize}
\raggedright
  \item \textquotedblleft{}\hi{मैं} \hi{हिंदी} \hi{बोलता} \hi{था}\textquotedblright{}
  \item the same sentence as a woman would say it
  \item which piece carries the habit and which carries the time
  \item \emph{I did not use to speak Hindi}, putting \hi{नहीं} in its usual seat
\end{itemize}

\begin{tcolorbox}[breakable,colback=teal!4,colframe=teal!35!black,title={Before you move on}]
How many words change between \emph{\hi{बोलता} \hi{हूँ}} and \emph{\hi{बोलता} \hi{था}}? (One.) What does the \textbf{-\hi{ता}} ending tell you, in either tense? (That the action is a habit rather than a single event.)
\end{tcolorbox}
\section[atīt kī tālikā]{atīt kī tālikā — the past, on one page}

\label{lesson:HI-C87-repaso-atit}



\begin{quote}
Name all three past endings and say who each one belongs to, before you look at the table.
\end{quote}

\begin{grammarlens}[title={Grammar Lens: the three endings}]
\noindent
\begin{tabularx}{\linewidth}{@{}>{\raggedright\arraybackslash}X>{\raggedright\arraybackslash}X@{}}
\toprule
\textbf{subject} & \textbf{earlier} \\
\midrule
\hi{मैं}, \hi{वह} — a man & \hi{था} \\
\hi{मैं}, \hi{वह} — a woman & \hi{थी} \\
\hi{हम}, \hi{वे}, \hi{आप}, \hi{तुम} & \hi{थे} \\
\bottomrule
\end{tabularx}

Three words, and their endings are \textbf{-\hi{आ}}, \textbf{-\hi{ई}} and \textbf{-\hi{ए}} — the same three you have been bending adjectives with since the colours.
\end{grammarlens}

\begin{grammarlens}[title={Grammar Lens: the swap, in every sentence you own}]
\noindent
\begin{tabularx}{\linewidth}{@{}>{\raggedright\arraybackslash}X>{\raggedright\arraybackslash}X@{}}
\toprule
\textbf{now} & \textbf{earlier} \\
\midrule
\hi{मैं} \hi{घर} \hi{पर} \hi{हूँ।} & \hi{मैं} \hi{घर} \hi{पर} \hi{था।} \\
\hi{वह} \hi{स्टेशन} \hi{पर} \hi{है।} & \hi{वह} \hi{स्टेशन} \hi{पर} \hi{था।} \\
\hi{मैं} \hi{हिंदी} \hi{बोलता} \hi{हूँ।} & \hi{मैं} \hi{हिंदी} \hi{बोलता} \hi{था।} \\
\bottomrule
\end{tabularx}

One word out, one word in. The place phrase, the postposition \textbf{\hi{पर}} and the word order all sit still while the tense moves around them.
\end{grammarlens}

\begin{grammarlens}[title={Grammar Lens: with \hi{नहीं}}]
\noindent
\begin{tabularx}{\linewidth}{@{}>{\raggedright\arraybackslash}X>{\raggedright\arraybackslash}X@{}}
\toprule
\textbf{sentence} & \textbf{meaning} \\
\midrule
\hi{हम} \hi{घर} \hi{पर} \hi{नहीं} \hi{थे।} & We were not at home. \\
\hi{वे} \hi{स्टेशन} \hi{पर} \hi{नहीं} \hi{थे।} & They were not at the station. \\
\bottomrule
\end{tabularx}
\end{grammarlens}

\subsection*{Your turn}
Say these aloud:

\begin{itemize}
\raggedright
  \item the three past endings in order
  \item \textquotedblleft{}\hi{मैं} \hi{घर} \hi{पर} \hi{था}\textquotedblright{} and then the same thing as a woman
  \item \textquotedblleft{}\hi{आप} \hi{कहाँ} \hi{थे}?\textquotedblright{}
  \item \textquotedblleft{}\hi{मैं} \hi{हिंदी} \hi{बोलता} \hi{था}\textquotedblright{}
  \item any one of those four with \hi{नहीं} added in the right seat
\end{itemize}

\begin{tcolorbox}[breakable,colback=teal!4,colframe=teal!35!black,title={Before you move on}]
Which past word does \hi{आप} take? (\textbf{\hi{थे}}.) Which side of the place does \hi{पर} go? (After it.) Which piece of \emph{\hi{बोलता} \hi{था}} carries the habit? (\textbf{-\hi{ता}}.)
\end{tcolorbox}
\section[āp kahā̃ the?]{\hi{आप} \hi{कहाँ} \hi{थे}? — the paper's own sentence, one step into the past}

\label{lesson:HI-C87-sintesis-atit}



\begin{quote}
Read this and say what it means before reading on.

\begin{quote}
\hi{मैं} \hi{घर} \hi{पर} \hi{हूँ।}
\end{quote}
\end{quote}

\begin{grammarlens}[title={Grammar Lens: the sentence a paper prints}]
That line is reading item 14 of the first practice paper, word for word. Until this chapter it had one unreadable piece in it, and the piece was \textbf{\hi{पर}}.

Now take one step:

\noindent
\begin{tabularx}{\linewidth}{@{}>{\raggedright\arraybackslash}X>{\raggedright\arraybackslash}X@{}}
\toprule
\textbf{the paper's line} & \textbf{one step back} \\
\midrule
\hi{मैं} \hi{घर} \hi{पर} \hi{हूँ।} & \hi{मैं} \hi{घर} \hi{पर} \hi{था।} \\
\bottomrule
\end{tabularx}

Everything that made the first line readable makes the second line readable too.
\end{grammarlens}

\subsection*{The exchange}
\begin{quote}
— \hi{आप} \hi{कहाँ} \hi{थे}? — \hi{मैं} \hi{घर} \hi{पर} \hi{था।}
\end{quote}

Three things are doing their job at once here. \textbf{\hi{कहाँ}} asks for a place. \textbf{\hi{पर}} attaches the answer to one. \textbf{\hi{थे}} is the past word \hi{आप} takes, and \textbf{\hi{था}} is the one the answer needs, because the answer is about one man speaking for himself.

A woman answering the same question says \textbf{\hi{मैं} \hi{घर} \hi{पर} \hi{थी}}, and the question itself does not change: the asker does not know yet who will answer.

\subsection*{How to answer — three ways}
Answer the question three ways, out loud:

\begin{itemize}
\raggedright
  \item you were at home
  \item you were at the station
  \item you were not at home
\end{itemize}

Then ask it of somebody you address as \hi{आप}, and listen for which of \hi{था}, \hi{थी} or \hi{थे} comes back.

\subsection*{Your turn}
Say these aloud:

\begin{itemize}
\raggedright
  \item \textquotedblleft{}\hi{आप} \hi{कहाँ} \hi{थे}?\textquotedblright{}
  \item \textquotedblleft{}\hi{मैं} \hi{घर} \hi{पर} \hi{था}\textquotedblright{}
  \item \textquotedblleft{}\hi{मैं} \hi{स्टेशन} \hi{पर} \hi{नहीं} \hi{था}\textquotedblright{}
  \item the answer a woman would give to the same question
\end{itemize}

\begin{tcolorbox}[breakable,colback=teal!4,colframe=teal!35!black,title={Before you move on}]
Which past word does the question to \hi{आप} use? (\textbf{\hi{थे}}.) Which one does a man use to answer about himself? (\textbf{\hi{था}}.) Which word in the answer puts him at a place? (\textbf{\hi{पर}}.)
\end{tcolorbox}
